% 埃尔德什·帕尔（Paul Erdős）（综述）
% license CCBYSA3
% type Wiki

本文根据 CC-BY-SA 协议转载翻译自维基百科\href{https://en.wikipedia.org/wiki/Paul_Erd\%C5\%91s}{相关文章}。

保罗·埃尔德什（英语：Paul Erdős，匈牙利语：Erdős Pál [ˈɛrdøːʃ ˈpaːl]；1913 年 3 月 26 日—1996 年 9 月 20 日）是一位匈牙利数学家。他是 20 世纪最具产出力的数学家之一，也是提出最多数学猜想的人之一。[2][3] 埃尔德什在离散数学、图论、数论、数学分析、逼近论、集合论以及概率论等领域追求并提出了大量问题。[4] 他的大部分研究集中在离散数学领域，解决了其中许多长期未解的难题。他大力倡导并作出重要贡献的拉姆齐理论，研究的是在何种条件下秩序必然出现。总体而言，埃尔德什的工作更倾向于解决已存在的开放问题，而非开辟或探索全新的数学领域。他一生共发表约 1500 篇数学论文，这一数字至今无人能及。[5]

他以“社会化数学”的实践方式以及古怪的生活方式闻名；《时代》杂志称他为“怪人中的怪人”。[6] 他坚信数学是一项社会性活动，因此过着漂泊不定的生活，唯一目的就是与其他数学家共同撰写论文。即使在晚年，他仍将清醒的每个时刻都献给数学研究；1996 年，他在华沙参加数学会议时去世。[7]

埃尔德什与合作者们的大量共同成果，促成了“埃尔德什数”的诞生——它表示某位数学家与埃尔德什之间的最短合著路径中的“步数”。
